% CERP Exam Question Bank - Set 2 (Questions 201-400)
% Additional 200 Unique Questions - Not Repeated from Set 1
% Compiled from ABA CERP Content Outline
% Compiled on: 2026-06-13

\documentclass[12pt, letterpaper]{article}
\usepackage[utf8]{inputenc}
\usepackage{geometry}
\usepackage{amsmath}
\usepackage{amsfonts}
\usepackage{amssymb}
\usepackage{titlesec}
\usepackage{enumitem}
\usepackage{xcolor}
\usepackage{fancyhdr}
\usepackage{multicol}

\geometry{margin=1in}
\setlength{\parindent}{0pt}
\setlength{\parskip}{0.8ex}

\pagestyle{fancy}
\fancyhf{}
\fancyhead[L]{\textbf{CERP Exam Prep - Set 2}}
\fancyhead[R]{\thepage}
\fancyfoot[C]{© 2026 Enterprise Risk Training}

\newcounter{questioncounter}
\setcounter{questioncounter}{201}

\newcommand{\question}{
    \par\noindent
    \textbf{\arabic{questioncounter}.}\quad
    \stepcounter{questioncounter}
}

\newcommand{\answer}[1]{
    \par\noindent
    \textbf{Answer:} #1
    \vspace{0.3cm}
    \par
}

\titleformat{\section}{\large\bfseries}{\thesection}{1em}{}
\titleformat{\subsection}{\normalsize\bfseries}{\thesubsection}{1em}{}

\begin{document}

\begin{center}
    \vspace*{1cm}
    {\huge \textbf{Certified Enterprise Risk Professional (CERP)}} \\[0.5cm]
    {\huge \textbf{Exam Question Bank - Set 2}} \\[1cm]
    {\Large 200 Additional Multiple-Choice Questions (201-400)} \\[1.5cm]
    {\large Based on ABA CERP Exam Content Outline} \\[0.3cm]
    {\large Domains 1-8: Risk Governance \& Risk Management} \\[2cm]
    {\large \textbf{Questions 201 - 400}} \\[0.3cm]
    \vfill
    {\large \today}
\end{center}

\newpage

% ============================================================
% DOMAIN 1: Risk Governance (Questions 201-216)
% ============================================================
\section{RISK GOVERNANCE}

\subsection{Domain 1: Board and Senior Management Oversight (Questions 201-216)}

\question Which committee is typically responsible for approving the annual risk appetite statement?
A) IT Committee
B) Compensation Committee
C) Board Risk Committee
D) Marketing Committee
\answer{C}

\question What is the primary difference between risk appetite and risk tolerance?
A) They are the same
B) Risk appetite is broad/strategic; risk tolerance is specific/quantitative limits
C) Risk tolerance is broader
D) Only risk appetite is regulatory
\answer{B}

\question Which behavior indicates a strong risk culture?
A) Hiding losses from management
B) Escalating issues promptly regardless of personal impact
C) Ignoring policy for speed
D) Blaming others for failures
\answer{B}

\question Senior management's role in risk governance includes:
A) Only reviewing audit reports
B) Implementing the risk appetite framework and ensuring adherence
C) Conducting external audits
D) Setting deposit rates
\answer{B}

\question What is "risk appetite" in the context of strategic planning?
A) A limit on employee bonuses
B) The amount and type of risk the bank pursues to achieve its strategy
C) The bank's marketing budget
D) The bank's office locations
\answer{B}

\question How often should the Board review the risk appetite statement?
A) Once every 5 years
B) At least annually and when circumstances change
C) Never
D) Only during a crisis
\answer{B}

\question Which is an example of poor risk culture?
A) Encouraging open discussion of risks
B) Retaliating against employees who raise concerns
C) Providing risk training
D) Having clear escalation policies
\answer{B}

\question What does "risk transparency" mean for the Board?
A) Only positive risks are reported
B) Full visibility into material risks, including emerging ones
C) Risks are hidden from regulators
D) Only losses are reported
\answer{B}

\question The three lines of defense model places independent assurance in which line?
A) First line (operations)
B) Second line (risk/compliance)
C) Third line (internal audit)
D) Fourth line (external consultants)
\answer{C}

\question What is a key responsibility of the Board's risk committee?
A) Approving individual loans
B) Reviewing risk management effectiveness and significant risk exposures
C) Hiring all staff
D) Setting daily trading limits
\answer{B}

\question Which action best demonstrates "tone at the middle"?
A) CEO memo
B) Front-line managers modeling risk-aware decision making
C) Board resolution
D) Regulatory fine
\answer{B}

\question What is a leading indicator of risk culture deterioration?
A) Increased training attendance
B) High voluntary risk reporting
C) Increased pressure to meet targets regardless of risk
D) Low employee turnover
\answer{C}

\question How does risk appetite link to capital planning?
A) It has no link
B) Risk appetite determines how much capital is needed to support risk-taking
C) Capital planning ignores risk
D) Only liquidity matters
\answer{B}

\question What is "credible challenge" in board meetings?
A) Always agreeing with management
B) Thoughtful questioning of management's risk assessments
C) Avoiding all conflict
D) Only challenging external parties
\answer{B}

\question Which is a sign of effective board risk oversight?
A) Board meets only once per year
B) Board receives dashboards with key risk indicators and exceptions
C) Board delegates all risk to one person
D) Board never asks questions
\answer{B}

\question What is the consequence of a weak risk culture?
A) Higher profitability
B) Undetected risk-taking leading to potential losses
C) Lower regulatory scrutiny
D) Reduced need for policies
\answer{B}

% ============================================================
% DOMAIN 2: Policies, Procedures and Limits (Questions 217-240)
% ============================================================

\subsection{Domain 2: Policies, Procedures and Limits (Questions 217-240)}

\question What is a "policy limit" in the context of risk management?
A) A suggestion
B) A boundary that defines acceptable risk-taking
C) A marketing target
D) An employee benefit
\answer{B}

\question Who should approve a new risk management policy?
A) Any employee
B) Appropriate level of management or Board based on policy significance
C) External auditor
D) Customer advisory board
\answer{B}

\question What is a typical limit for liquidity risk?
A) Loan-to-value ratio
B) Minimum liquidity coverage ratio (LCR)
C) Employee turnover rate
D) Customer satisfaction score
\answer{B}

\question A policy breach that is not escalated may result in:
A) No consequences
B) Accumulated unrecognized risk and regulatory sanctions
C) Automatic bonus
D) Reduced workload
\answer{B}

\question What should be included in a policy exception request?
A) Only the amount
B) Business justification, risk assessment, and proposed approver
C) Employee's favorite color
D) Customer's birth date
\answer{B}

\question How often should operational risk limits be reviewed?
A) Never
B) At least annually and following material changes
C) Only after a loss
D) Only when regulators ask
\answer{B}

\question Which is a regulatory expectation for policy governance?
A) Policies can be verbal
B) Policies must be written, approved, and accessible
C) Policies need no version control
D) Policies can contradict each other
\answer{B}

\question What is the purpose of a risk limit escalation process?
A) To ignore breaches
B) To ensure breaches are reported to appropriate management for action
C) To automatically approve all breaches
D) To hide breaches from the Board
\answer{B}

\question A concentration limit on a single counterparty is designed to prevent:
A) Diversification
B) Excessive loss from a single default
C) Low returns
D) Operational efficiency
\answer{B}

\question What is the difference between a policy and a procedure?
A) No difference
B) Policy states "what" and "why"; procedure states "how"
C) Procedure is more important
D) Policy is verbal
\answer{B}

\question Which is an example of a risk limit based on capital?
A) \$10,000 expense limit
B) Maximum loss not to exceed 5\% of Tier 1 capital
C) 10 employee limit
D) 100 customer limit
\answer{B}

\question What is a "hard limit" vs. a "soft limit"?
A) Soft limit cannot be exceeded; hard limit can with approval
B) Hard limit cannot be exceeded; soft limit may be exceeded with approval
C) Both are the same
D) Neither requires approval
\answer{B}

\question Who is responsible for monitoring compliance with policy limits?
A) Only internal audit
B) First line (business) with second line oversight
C) Only the Board
D) External regulators
\answer{B}

\question A policy exception that is not properly documented may be treated as:
A) Approved
B) A policy violation
C) Ignored
D) A best practice
\answer{B}

\question What is a typical limit for credit risk concentration?
A) Unlimited to one borrower
B) Maximum percentage of capital to a single borrower
C) No limit on industry exposure
D) Unlimited unsecured lending
\answer{B}

\question What is the role of the Board in approving policy exceptions?
A) Approve every exception
B) Approve material or significant exceptions above delegated authority
C) No role
D) Only review after the fact
\answer{B}

\question A breach of a risk limit should trigger:
A) No action
B) Investigation, escalation, and remediation plan
C) Immediate termination of all staff
D) Ignoring the breach
\answer{B}

\question What is a "dispensation" in risk management?
A) A permanent policy change
B) A time-limited permission to operate outside a policy limit
C) An audit finding
D) A regulatory fine
\answer{B}

\question Policy limits should be documented in:
A) Personal emails
B) The risk appetite framework or policy documents
C) Verbal agreements
D) Marketing brochures
\answer{B}

\question Which is a key element of policy governance?
A) No review process
B) Scheduled periodic review and version control
C) Unlimited exceptions without approval
D) No approval requirements
\answer{B}

\question A counterparty credit limit is an example of:
A) Operational risk limit
B) Credit risk limit
C) Strategic risk limit
D) Reputational risk limit
\answer{B}

\question What should happen when a policy limit is breached?
A) Nothing
B) Immediate notification to risk management and remediation
C) Only fix if material
D) Wait until month-end
\answer{B}

\question A risk limit based on earnings might state:
A) No limit
B) Maximum loss of X\% of net income
C) Unlimited losses
D) No earnings constraint
\answer{B}

\question Who typically sets the overall risk limit framework?
A) Front-line staff
B) Board of Directors or Board-level committee
C) External vendors
D) Customers
\answer{B}

% ============================================================
% DOMAIN 3: MIS and Data Governance (Questions 241-255)
% ============================================================

\subsection{Domain 3: Management Information Systems (Questions 241-255)}

\question What is the primary risk of manual reporting?
A) Too fast
B) Higher risk of human error and manipulation
C) No risk
D) Always accurate
\answer{B}

\question Which is a key component of data lineage?
A) Deleting old data
B) Tracking data from source to report
C) Ignoring data sources
D) Only using external data
\answer{B}

\question What is a "data dictionary" used for?
A) Storing actual data
B) Defining data fields, formats, and business meanings
C) Deleting data
D) Marketing purposes
\answer{B}

\question Which is an example of a data quality issue?
A) Accurate data
B) Incomplete or duplicate data
C) Timely data
D) Relevant data
\answer{B}

\question What is the purpose of a data governance committee?
A) To reduce data quality
B) To oversee data policies, quality, and access
C) To delete all data
D) To limit data access to one person
\answer{B}

\question Which is a control for ensuring data completeness?
A) No validation
B) Reconciliation between source systems and reports
C) Allowing any user to modify data
D) No audit trail
\answer{B}

\question What is a "model validation"?
A) Using a model without review
B) Independent review of model assumptions, inputs, and outputs
C) Deleting the model
D) Marketing the model
\answer{B}

\question An MIS report for the Board should focus on:
A) Every transaction
B) Aggregated risk exposures, limit exceptions, and trends
C) Employee daily tasks
D) Office supplies
\answer{B}

\question What is a key risk of system access restrictions?
A) Too many people have access
B) Authorized users cannot obtain necessary data in a timely manner
C) No risk
D) Data is too secure
\answer{B}

\question Which is a best practice for risk data aggregation?
A) Siloed data in each department
B) Centralized or integrated data warehouse with defined data models
C) No data standards
D) Manual spreadsheets only
\answer{B}

\question What is the purpose of "user acceptance testing" (UAT) for a risk system?
A) To skip testing
B) To ensure the system meets business requirements before go-live
C) To delete the system
D) To market the system
\answer{B}

\question A data quality KPI might track:
A) Employee satisfaction
B) Percentage of data fields populated and error rates
C) Marketing spend
D) Office temperature
\answer{B}

\question What is a "data owner" responsible for?
A) Nothing
B) Data quality, access, and definition for their domain
C) Only data deletion
D) Only data storage
\answer{B}

\question Which is an example of a critical data element for credit risk?
A) Employee name
B) Borrower's loan balance and payment status
C) Office address
D) Marketing campaign name
\answer{B}

\question What is the risk of poor data governance?
A) Perfect reporting
B) Misinformed decisions, regulatory fines, and financial loss
C) Lower costs
D) Faster reporting
\answer{B}

% ============================================================
% DOMAIN 4: Control Framework (Questions 256-270)
% ============================================================

\subsection{Domain 4: Control Framework (Questions 256-270)}

\question What is a "compensating control"?
A) A control that does nothing
B) An alternative control used when a primary control is not feasible
C) A control that increases risk
D) A marketing control
\answer{B}

\question Which is an example of a corrective control?
A) Pre-approval
B) Disaster recovery plan following an incident
C) Segregation of duties
D) Data validation
\answer{B}

\question What is "control effectiveness"?
A) Whether a control exists
B) Whether a control is designed properly and operates as intended
C) The cost of the control
D) The number of controls
\answer{B}

\question Which is a key principle of COSO's control environment?
A) No ethical values
B) Commitment to integrity and ethical values
C) No oversight
D) No accountability
\answer{B}

\question What is the difference between "preventive" and "detective" controls?
A) No difference
B) Preventive stops errors; detective finds errors after occurrence
C) Detective stops errors
D) Preventive finds errors
\answer{B}

\question What is a "control gap"?
A) A control that works perfectly
B) A risk not addressed by any control
C) Too many controls
D) An automated control
\answer{B}

\question Which role performs independent testing of controls?
A) First line
B) Second line
C) Third line (internal audit)
D) Front-line staff
\answer{C}

\question What is "control automation" intended to reduce?
A) Accuracy
B) Human error and processing delays
C) Speed
D) Compliance
\answer{B}

\question Which is a control for fraud risk?
A) No segregation of duties
B) Mandatory vacations and job rotation
C) Unlimited access to all systems
D) No approval requirements
\answer{B}

\question What is a "key control"?
A) An unimportant control
B) A critical control that is essential to mitigating a material risk
C) A control that is never tested
D) A control without documentation
\answer{B}

\question How often should controls be tested for operating effectiveness?
A) Never
B) Periodically based on risk (e.g., annually or more frequently)
C) Only when there is a loss
D) Once every 10 years
\answer{B}

\question What is the purpose of a "control self-assessment" (CSA)?
A) To eliminate controls
B) To allow process owners to evaluate their own control effectiveness
C) To outsource controls
D) To ignore risks
\answer{B}

\question Which is a control for regulatory compliance?
A) Ignoring regulations
B) Compliance monitoring and testing program
C) No training
D) No policies
\answer{B}

\question What is "segregation of duties" intended to prevent?
A) Efficiency
B) Fraud and error by ensuring no single person has excessive control
C) Teamwork
D) Automation
\answer{B}

\question A detective control that identifies a fraud after it occurs:
A) Prevents the fraud
B) Allows timely remediation and future prevention
C) Ignores the fraud
D) Rewards the fraud
\answer{B}

% ============================================================
% DOMAIN 5: Risk Identification (Questions 271-295)
% ============================================================

\subsection{Domain 5: Risk Identification (Questions 271-295)}

\question Which is a compliance risk event?
A) Loan default
B) Violation of anti-money laundering (AML) regulations
C) Interest rate change
D) System upgrade
\answer{B}

\question What is "emerging risk"?
A) A risk that has already materialized
B) A new or evolving risk that may become material in the future
C) A risk that never happens
D) A risk that is always decreasing
\answer{B}

\question Which is a third-party risk?
A) Employee error
B) Vendor data breach affecting the bank
C) Internal policy violation
D) Board decision
\answer{B}

\question What is a "risk event"?
A) A planned activity
B) An occurrence that causes actual or potential negative impact
C) A routine transaction
D) An employee meeting
\answer{B}

\question Which is a strategic risk?
A) Processing error
B) Poor business decision or failure to adapt to market changes
C) Loan default
D) System outage
\answer{B}

\question What is the purpose of a "risk register"?
A) To ignore risks
B) To document identified risks, their assessments, and ownership
C) To delete risks
D) To market products
\answer{B}

\question Which is an example of a reputational risk event?
A) Loan payment
B) Negative media coverage of unethical practices
C) Interest rate cut
D) New product launch
\answer{B}

\question What is a "silent risk"?
A) A well-known risk
B) A risk that is present but not identified or documented
C) A loud risk
D) A risk with no impact
\answer{B}

\question Which is a source of risk from a third-party vendor?
A) No risk
B) Financial instability of the vendor leading to service disruption
C) Perfect performance
D) Low fees
\answer{B}

\question What is a "risk taxonomy"?
A) A list of animals
B) A structured classification of risk categories and sub-categories
C) A marketing tool
D) A software system
\answer{B}

\question Which is an example of a model risk?
A) Using a model correctly
B) Using an incorrect or poorly validated model for decision-making
C) No model usage
D) Manual underwriting
\answer{B}

\question What is a "control weakness" in risk identification?
A) A strong control
B) A deficiency that increases the likelihood of a risk event
C) An automated control
D) A preventive control
\answer{B}

\question Which is an external risk factor?
A) Employee skills
B) Regulatory change
C) Internal policies
D) Bank culture
\answer{B}

\question What is a "risk indicator"?
A) A guarantee of safety
B) A metric that signals changes in risk level
C) A final loss number
D) A marketing KPI
\answer{B}

\question Which is a legal risk event?
A) Loan origination
B) Lawsuit alleging breach of contract
C) Customer deposit
D) Branch opening
\answer{B}

\question What is the value of a "risk and control matrix" (RACM)?
A) No value
B) Maps risks to controls and identifies gaps
C) Deletes all controls
D) Only for IT
\answer{B}

\question Which is an internal risk factor?
A) Economic downturn
B) Inadequate internal training and skills gaps
C) New regulation
D) Competitor action
\answer{B}

\question What is the first step in risk identification?
A) Ignore all inputs
B) Understanding business objectives and processes
C) Buying software
D) Hiring consultants
\answer{B}

\question Which is a "horizontal risk" that affects multiple business lines?
A) A risk isolated to one product
B) Cyber risk affecting the entire enterprise
C) A risk with no impact
D) A risk that is always acceptable
\answer{B}

\question What is a "risk trigger"?
A) A control
B) An event or condition that causes a risk to materialize
C) A risk response
D) A risk appetite statement
\answer{B}

\question Which is a geopolitical risk?
A) Loan default
B) Trade sanctions or political instability in a country of operation
C) System failure
D) Employee fraud
\answer{B}

\question What is "risk awareness"?
A) Ignoring risks
B) Understanding that risks exist and can impact objectives
C) Only looking at past losses
D) Avoiding all training
\answer{B}

\question Which is a risk from rapid business growth?
A) No risk
B) Strained controls and operational capacity
C) Higher profitability with no downsides
D) Lower risk
\answer{B}

\question What is a "risk universe"?
A) Outer space
B) The complete set of all possible risks that could affect the organization
C) Only financial risks
D) Only operational risks
\answer{B}

\question Which is a risk from new product introduction?
A) Guaranteed success
B) Inadequate risk assessment before launch
C) No regulatory review
D) Automatic profitability
\answer{B}

% ============================================================
% DOMAIN 6: Risk Measurement (Questions 296-320)
% ============================================================

\subsection{Domain 6: Risk Measurement and Evaluation (Questions 296-320)}

\question What is "expected loss" (EL) in credit risk?
A) Loss that never happens
B) Probability of Default (PD) × Exposure at Default (EAD) × Loss Given Default (LGD)
C) Maximum possible loss
D) Regulatory capital only
\answer{B}

\question What is "unexpected loss" (UL)?
A) Same as expected loss
B) The volatility of loss around expected loss (requires capital)
C) A loss that is always zero
D) Only operational loss
\answer{B}

\question Which is a measure of market risk?
A) Loan-to-value
B) Value at Risk (VaR) or Duration
C) Employee turnover
D) Customer complaints
\answer{B}

\question What is "stress testing"?
A) A type of exercise
B) Evaluating the impact of severe but plausible scenarios
C) Regular business as usual
D) Ignoring tail risks
\answer{B}

\question What is "reverse stress testing"?
A) Testing normal conditions
B) Identifying scenarios that would cause business failure
C) No testing
D) Only regulatory testing
\answer{B}

\question Which is a key indicator for interest rate risk?
A) Loan volume
B) Earnings at Risk (EaR) or Economic Value of Equity (EVE) sensitivity
C) Number of employees
D) Marketing spend
\answer{B}

\question What is "risk correlation"?
A) Risks are independent
B) How different risks move together (can amplify or offset)
C) No relationship between risks
D) Only one risk matters
\answer{B}

\question Which is a measure of operational risk?
A) Stock price
B) Historical loss data and scenario analysis
C) Loan default rate
D) Interest rate curve
\answer{B}

\question What is "risk appetite tolerance"?
A) No limit
B) Specific quantitative boundaries for each risk
C) Only qualitative statements
D) The same as risk appetite
\answer{B}

\question Which is a liquidity risk metric?
A) Loan-to-deposit ratio and High-Quality Liquid Assets (HQLA)
B) Employee satisfaction
C) Number of ATMs
D) Marketing ROI
\answer{A}

\question What is "sensitivity analysis"?
A) No analysis
B) Changing one input at a time to see impact on output
C) Changing everything at once
D) Only qualitative analysis
\answer{B}

\question Which is a key risk indicator for model risk?
A) Number of models without validation
B) Number of employees
C) Branch locations
D) Marketing campaigns
\answer{A}

\question What is "backtesting" in risk measurement?
A) Testing a model on historical data to compare predictions to actual outcomes
B) Testing only once
C) No testing
D) Testing only future data
\answer{A}

\question What is "concentration risk" measured by?
A) Herfindahl-Hirschman Index (HHI) or top borrower percentage
B) Number of employees
C) Office square footage
D) Marketing budget
\answer{A}

\question What is "probability of default" (PD)?
A) The loss if default occurs
B) The likelihood that a borrower will default within a given time horizon
C) The exposure amount
D) The recovery rate
\answer{B}

\question What is "loss given default" (LGD)?
A) Probability of default
B) The percentage of exposure lost if default occurs
C) The total exposure
D) The recovery amount
\answer{B}

\question Which is a KRIs for compliance risk?
A) Number of customer complaints
B) Number of regulatory violations or self-identified issues
C) Loan growth rate
D) Deposit balance
\answer{B}

\question What is "risk velocity" measurement?
A) Measuring size
B) Assessing how quickly a risk can materialize and escalate
C) Measuring color
D) Measuring location
\answer{B}

\question Which is a key indicator for cyber risk?
A) Number of desktops
B) Number of unpatched critical systems or security incidents
C) Number of employees
D) Office location
\answer{B}

\question What is "risk aggregation"?
A) Looking at risks separately
B) Combining risk exposures to see total enterprise risk
C) Ignoring correlations
D) Only looking at one risk
\answer{B}

\question What is a "risk limit utilization report"?
A) A marketing report
B) Shows current usage versus approved risk limits
C) An HR report
D) A facilities report
\answer{B}

\question Which is a measure of reputational risk?
A) Loan-to-value
B) Net Promoter Score (NPS) or media sentiment analysis
C) Interest rate duration
D) Stock price only
\answer{B}

\question What is "economic capital"?
A) Regulatory capital only
B) The amount of capital needed to cover unexpected losses at a confidence level
C) Book equity
D) Debt
\answer{B}

\question What is "risk-adjusted return on capital" (RAROC)?
A) Raw return
B) Return divided by economic capital (risk-adjusted profitability)
C) Only cost
D) Only revenue
\answer{B}

\question What is "tail risk"?
A) Common risk
B) Risk of extreme outcomes beyond normal distribution assumptions
C) No risk
D) Certain risk
\answer{B}

% ============================================================
% DOMAIN 7: Risk Responses (Questions 321-345)
% ============================================================

\subsection{Domain 7: Risk Responses (Questions 321-345)}

\question When is "risk avoidance" the appropriate response?
A) When risk is low
B) When the risk exceeds the potential reward and cannot be mitigated effectively
C) Always
D) Never
\answer{B}

\question What is "risk transfer"?
A) Keeping all risk
B) Shifting risk to a third party (e.g., insurance, hedging)
C) Ignoring risk
D) Accepting risk
\answer{B}

\question What is "root cause analysis"?
A) Ignoring causes
B) A structured method to identify the underlying source of an issue
C) Only looking at symptoms
D) A marketing technique
\answer{B}

\question Which is an example of risk mitigation?
A) Doing nothing
B) Implementing controls to reduce likelihood or impact
C) Avoiding all activity
D) Buying insurance only
\answer{B}

\question What is "risk acceptance"?
A) Ignoring risk without decision
B) Acknowledging risk and deciding not to take action (within appetite)
C) Always mitigating
D) Always transferring
\answer{B}

\question What is an "issue" in risk management?
A) A successful process
B) A identified deficiency or failure in controls or compliance
C) A routine report
D) A training session
\answer{B}

\question What is an "action plan" for an issue?
A) No plan
B) A documented set of steps to remediate the issue with timelines and ownership
C) A marketing brochure
D) An employee evaluation
\answer{B}

\question What is "issue validation"?
A) Assuming issue is fixed
B) Independent confirmation that remediation was effective
C) Ignoring closure
D) Closing without evidence
\answer{B}

\question Which is a risk response for a high-severity, high-likelihood risk?
A) Ignore
B) Aggressive mitigation or avoidance
C) Accept
D) Defer indefinitely
\answer{B}

\question What is "residual risk" after risk transfer?
A) The same as inherent risk
B) The risk remaining after insurance or hedging
C) Zero risk
D) Higher risk
\answer{B}

\question What is the "5 Whys" technique used for?
A) Marketing
B) Root cause analysis by repeatedly asking "why"
C) Financial reporting
D) HR evaluations
\answer{B}

\question Which is a detective mitigation activity?
A) Pre-approval of transactions
B) Reconciliation and exception reporting
C) Segregation of duties
D) Physical security
\answer{B}

\question What is an "issue severity rating"?
A) Ignoring issues
B) Classification of issue impact (e.g., high, medium, low)
C) The issue owner's name
D) The issue date
\answer{B}

\question When is risk transfer most effective?
A) For highly predictable, small losses
B) For low-frequency, high-severity risks (e.g., catastrophes)
C) For all risks
D) For no risks
\answer{B}

\question What is a "corrective action plan" (CAP)?
A) No action
B) A plan to fix a regulatory or audit finding
C) A marketing plan
D) A budget plan
\answer{B}

\question Which is an example of risk acceptance?
A) Buying insurance
B) Self-insuring within appetite and setting aside reserves
C) Terminating a product line
D) Adding controls
\answer{B}

\question What is "issue tracking"?
A) Ignoring issues
B) A system to log, monitor, and report issues through closure
C) Deleting issues
D) Only reporting positive items
\answer{B}

\question What is the role of the second line in issue management?
A) Own the issue
B) Oversee, challenge, and validate remediation
C) Ignore issues
D) Only originate issues
\answer{B}

\question Which is a preventive mitigation activity?
A) Detective review after transaction
B) Approval authority before transaction execution
C) Reconciliation after loss
D) Investigation after event
\answer{B}

\question What is "risk dispensation"?
A) Permanent elimination of risk
B) Temporary permission to operate outside a limit with conditions
C) Ignoring all limits
D) Automatic approval
\answer{B}

\question Which is a risk response for a low-impact, low-likelihood risk?
A) Intensive mitigation
B) Acceptance or monitoring only
C) Avoidance
D) Transfer
\answer{B}

\question What is the purpose of a "risk response plan"?
A) To ignore risks
B) To document actions to address identified risks
C) To create more risk
D) To eliminate all planning
\answer{B}

\question Which is an example of risk avoidance?
A) Hedging with derivatives
B) Discontinuing a high-risk product line entirely
C) Purchasing insurance
D) Accepting the risk
\answer{B}

\question What is "inherent risk" in risk response evaluation?
A) Risk after response
B) Risk before any response is applied
C) Residual risk
D) Transfer risk
\answer{B}

\question What is the goal of risk mitigation?
A) Eliminate all risk (impossible)
B) Reduce risk to within appetite using cost-effective controls
C) Ignore risk
D) Increase risk
\answer{B}

% ============================================================
% DOMAIN 8: Risk Monitoring (Questions 346-370)
% ============================================================

\subsection{Domain 8: Risk Monitoring (Questions 346-370)}

\question What is a "key risk indicator" (KRI) trigger?
A) A normal operating level
B) A pre-defined threshold that signals need for action
C) A final loss
D) A marketing event
\answer{B}

\question What is a "risk dashboard"?
A) A car part
B) A visual summary of key risk metrics, limit utilization, and exceptions
C) A type of computer
D) A marketing flyer
\answer{B}

\question Which is a credit risk KRI?
A) System uptime
B) 30+ day delinquency rate
C) Employee turnover
D) Number of branches
\answer{B}

\question What is "control monitoring"?
A) Ignoring controls
B) Ongoing assessment of whether controls are operating effectively
C) Deleting controls
D) Only reviewing annually
\answer{B}

\question Which is a financial KRI for earnings risk?
A) Operational loss count
B) Volatility of net interest margin (NIM)
C) Employee satisfaction
D) Number of policies
\answer{B}

\question What is "heat mapping" in risk reporting?
A) A weather report
B) A visual method to show risk levels by likelihood and impact
C) A cooking technique
D) An IT diagnostic
\answer{B}

\question Which is a non-financial KRI?
A) Loan balance
B) System downtime hours per month
C) Deposit amount
D) Capital ratio
\answer{B}

\question What is the appropriate escalation level for a breach of a critical risk limit?
A) No escalation
B) Immediate escalation to senior management and possibly Board
C) Escalate only at year-end
D) Escalate to marketing
\answer{B}

\question Which is a KRI for third-party risk?
A) Number of employees
B) Vendor cybersecurity assessment score or overdue findings
C) Customer satisfaction
D) Loan volume
\answer{B}

\question What is "ad hoc reporting"?
A) Scheduled daily report
B) Unscheduled report created to address a specific question or event
C) No reporting
D) Annual report only
\answer{B}

\question Which is a leading KRI?
A) Past losses
B) Number of overdue risk self-assessment findings
C) Final financial restatement
D) Regulatory fine amount
\answer{B}

\question What is a "lagging KRI"?
A) Future prediction
B) Measure of past events (e.g., actual losses)
C) Current exposure
D) Limit utilization
\answer{B}

\question Which is a KRI for compliance risk monitoring?
A) Loan growth
B) Percentage of completed required compliance trainingC) Deposit balance
D) Stock price
\answer{B}

\question What is the purpose of "risk reporting distribution"?
A) To hide information
B) To ensure the right people get the right information at the right time
C) To overwhelm everyone
D) To avoid action
\answer{B}

\question Which is a KRI for operational risk?
A) Loan default rate
B) Number of failed internal transactions or processing errors
C) Interest rate sensitivity
D) Credit spread
\answer{B}

\question What is "trend analysis" in risk monitoring?
A) Looking at one data point
B) Comparing risk metrics over time to identify patterns
C) Ignoring historical data
D) Only looking at current data
\answer{B}

\question Which is a KRI for liquidity risk monitoring?
A) Loan-to-value ratio
B) Borrowing capacity or available HQLA
C) Employee count
D) Number of ATMs
\answer{B}

\question What is a "risk appetite dashboard"?
A) A marketing dashboard
B) A report showing risk levels compared to appetite limits
C) An HR dashboard
D) A facilities dashboard
\answer{B}

\question Which is an example of automated control monitoring?
A) Manual review of every transaction
B) System-generated alerts for transactions exceeding limits
C) No monitoring
D) Annual audit only
\answer{B}

\question What is the role of the Board in risk monitoring?
A) Monitor every transaction
B) Receive regular reports on top risks, limit breaches, and emerging issues
C) Execute daily monitoring
D) No role
\answer{B}

\question Which is a KRI for strategic risk?
A) System downtime
B) Number of missed strategic milestones or budget variances
C) Loan default
D) Fraud loss
\answer{B}

\question What is "exception reporting"?
A) Reporting all transactions
B) Reporting only items that deviate from policy or limits
C) Reporting nothing
D) Reporting only favorable items
\answer{B}

\question Which is a KRI for model risk?
A) Number of models
B) Number of model limitations or overdue validations
C) Employee count
D) Branch count
\answer{B}

\question What is "continuous monitoring"?
A) Monitoring once per year
B) Ongoing, real-time or near-real-time risk surveillance
C) No monitoring
D) Manual monitoring only
\answer{B}

\question Which is a risk-based recommendation from monitoring?
A) Ignore all findings
B) Adjust limits, enhance controls, or change risk response
C) Delete all data
D) Hide results
\answer{B}

% ============================================================
% FINAL QUESTIONS 371-400 (Mixed Domains)
% ============================================================

\subsection{Comprehensive Review (Questions 371-400)}

\question Which regulation requires management to assess internal controls over financial reporting?
A) GDPR
B) Sarbanes-Oxley Act (SOX) Section 404
C) CCPA
D) HIPAA
\answer{B}

\question What is a "matter requiring attention" (MRA)?
A) A commendation
B) A regulatory finding requiring corrective action
C) A bonus approval
D) A marketing approval
\answer{B}

\question Which is a civil money penalty (CMP)?
A) A reward
B) A financial penalty imposed by regulators for violations
C) A tax deduction
D) An employee bonus
\answer{B}

\question What is the purpose of an internal control "narrative"?
A) A story
B) Written documentation of control design and operation
C) A marketing brochure
D) An employee handbook
\answer{B}

\question Which is an example of a compensating control?
A) No control
B) Manual review where automated control fails
C) Deleting all controls
D) Ignoring the risk
\answer{B}

\question What is "risk taxonomy" used for?
A) Categorizing animals
B) Standardizing risk language across the organization
C) Marketing segmentation
D) Employee classification
\answer{B}

\question Which is a key element of an RCSA?
A) Only financial risks
B) Process identification, risk assessment, control evaluation, action plans
C) No action plans
D) Only inherent risk
\answer{B}

\question What is "stress testing" for liquidity risk?
A) Testing employee stress
B) Assessing survival period under a liquidity freeze scenario
C) Testing marketing stress
D) Testing IT stress
\answer{B}

\question Which is an example of a key control?
A) An insignificant control
B) A control that prevents material misstatement or loss
C) A control that is never tested
D) A control without an owner
\answer{B}

\question What is the purpose of a "risk committee charter"?
A) To ignore risk
B) To document the committee's authority, responsibilities, and composition
C) To market products
D) To hire staff
\answer{B}

\question Which is a key indicator for fraud risk?
A) Loan growth
B) Number of policy override exceptions without justification
C) Deposit balance
D) Employee headcount
\answer{B}

\question What is "model governance"?
A) No oversight
B) Framework for model development, validation, use, and monitoring
C) Deleting models
D) Only using one model
\answer{B}

\question Which is a risk from regulatory change?
A) No risk
B) Increased compliance costs or business model changes
C) Guaranteed profit
D) Reduced risk
\answer{B}

\question What is "risk appetite statement"?
A) A marketing document
B) A written articulation of the types and amount of risk the organization is willing to accept
C) An HR policy
D) A facilities manual
\answer{B}

\question Which is a KRI for information security risk?
A) Loan volume
B) Number of unpatched critical vulnerabilities
C) Customer deposits
D) Branch locations
\answer{B}

\question What is the purpose of "escalation thresholds"?
A) To never escalate
B) To define when and to whom an issue must be raised
C) To hide issues
D) To delay action
\answer{B}

\question Which is an example of a detective control?
A) Pre-approval of expenses
B) Monthly bank reconciliation
C) Segregation of duties
D) Physical locks
\answer{B}

\question What is "control testing"?
A) Ignoring controls
B) Procedures to verify control design and operating effectiveness
C) Deleting controls
D) Only documenting controls
\answer{B}

\question Which is a risk from outsourcing?
A) No risk
B) Loss of control and potential service disruption
C) Guaranteed cost savings
D) No regulatory oversight
\answer{B}

\question What is the role of internal audit in risk management?
A) Own all risks
B) Provide independent assurance on risk management and control effectiveness
C) Execute risk-taking
D) Set risk appetite
\answer{B}

\question Which is a KRI for interest rate risk?
A) Loan growth rate
B) Duration gap or repricing mismatch
C) Employee turnover
D) Customer complaints
\answer{B}

\question What is a "risk assessment methodology"?
A) A random process
B) A structured approach to identify, analyze, and evaluate risks
C) A marketing plan
D) A hiring process
\answer{B}

\question Which is an example of a corrective control?
A) Pre-approval
B) Incident response plan following a breach
C) Segregation of duties
D) Data validation
\answer{B}

\question What is the purpose of a "risk limit breach report"?
A) To ignore breaches
B) To notify management and track remediation of limit exceedances
C) To hide breaches
D) To approve breaches automatically
\answer{B}

\question Which is a key element of the COSO framework?
A) Only monitoring
B) Information and communication
C) Only control activities
D) Only risk assessment
\answer{B}

\question What is "risk capacity"?
A) The same as risk appetite
B) The maximum amount of risk an organization can absorb without failing
C) A risk limit
D) A policy exception
\answer{B}

\question Which is an example of a preventive control?
A) Detective review
B) Physical security and access controls
C) Reconciliation after transaction
D) Investigation of loss
\answer{B}

\question What is a "risk register" used for?
A) Deleting risks
B) Tracking identified risks, assessments, and action owners
C) Marketing
D) HR purposes
\answer{B}

\question Which is a KRI for reputational risk?
A) Loan default rate
B) Negative media mentions or social media sentiment
C) Capital ratio
D) Interest rate
\answer{B}

\question What is the purpose of an "audit trail"?
A) To delete records
B) To document who did what and when for accountability
C) To hide transactions
D) To slow down processes
\answer{B}
\newpage

% ============================================================
% ANSWER KEY (Questions 201-400)
% ============================================================
\section*{Answer Key - Questions 201-400}
\begin{multicols}{4}
\small
\begin{enumerate}[start=201]
\item C
\item B
\item B
\item B
\item B
\item B
\item B
\item B
\item C
\item B
\item B
\item C
\item B
\item B
\item B
\item B
\item B
\item B
\item B
\item B
\item B
\item B
\item B
\item B
\item B
\item B
\item B
\item B
\item B
\item B
\item B
\item B
\item B
\item B
\item B
\item B
\item B
\item B
\item B
\item B
\item B
\item B
\item B
\item B
\item B
\item B
\item B
\item B
\item B
\item B
\item B
\item B
\item B
\item B
\item B
\item B
\item B
\item B
\item B
\item B
\item B
\item B
\item B
\item B
\item B
\item B
\item B
\item B
\item B
\item B
\item B
\item B
\item B
\item B
\item B
\item B
\item B
\item B
\item B
\item B
\item B
\item B
\item B
\item B
\item B
\item B
\item B
\item B
\item B
\item B
\item B
\item B
\item B
\item B
\item B
\item B
\item B
\item B
\item B
\item B
\item B
\item B
\item B
\item B
\item B
\item B
\item B
\item B
\item B
\item B
\item B
\item B
\item B
\item B
\item B
\item B
\item B
\item B
\item B
\item B
\item B
\item B
\item B
\item B
\item B
\item B
\item B
\item B
\item B
\item B
\item B
\item B
\item B
\item B
\item B
\item B
\item B
\item B
\item B
\item B
\item B
\item B
\item B
\item B
\item B
\item B
\item B
\item B
\item B
\item B
\item B
\item B
\item B
\item B
\item B
\item B
\item B
\item B
\item B
\item B
\item B
\item B
\item B
\item B
\item B
\item B
\item B
\item B
\item B
\item B
\item B
\item B
\item B
\item B
\item B
\item B
\item B
\item B
\item B
\item B
\item B
\item B
\item B
\item B
\item B
\item B
\item B
\item B
\item B
\item B
\item B
\item B
\item B
\item B
\item B
\item B
\item B
\item B
\item B
\item B
\end{enumerate}
\end{multicols}

\end{document}